% figures/dynamics/track_state_vs_vehicle_state.tex

\begin{tikzpicture}[
  node distance=1.15cm,
  box/.style={
    draw,
    rounded corners,
    align=center,
    minimum width=4.8cm,
    minimum height=0.9cm,
    font=\small
  },
  smallbox/.style={
    draw,
    rounded corners,
    align=center,
    minimum width=3.8cm,
    minimum height=0.8cm,
    font=\footnotesize
  },
  arrow/.style={->, thick}
]

\node[box] (track) {Track state\\\(\mathrm{pose3}(s)\)};
\node[box, right=3.0cm of track] (vehicle) {Vehicle state\\\(x_v(s,t)\)};

\node[smallbox, below=of track] (railway) {railway frame\\cant, profile, geometry};
\node[smallbox, below=of vehicle] (motion) {motion response\\speed, balance, comfort};

\draw[arrow] (track) -- node[above, font=\footnotesize] {runtime interpretation} (vehicle);
\draw[arrow] (track) -- (railway);
\draw[arrow] (vehicle) -- (motion);

\node[align=center, font=\small, below=2.7cm of $(track)!0.5!(vehicle)$]
{Track state and vehicle state are coupled,\\but they are not identical.};

\end{tikzpicture}
